% GENERATED by assemble_submissions.py -- do not edit.
% These blocks were lifted out of the body to meet ARR's 8-page
% content limit; appendices sit after the references and are exempt.

\section{The resolution floor of the randomization test}
\label{app:floor}

The argument is worth stating because the floor is a property of the design rather
than of the data. Our statistic is $|\mathrm{mean}(a) - \mathrm{mean}(b)|$. When
the two arms are the same size, the complement of any admissible group-$a$ subset
is itself an admissible group-$a$ subset, and swapping the groups negates the
difference while preserving its absolute value. Every regrouping at the observed
extreme therefore has a partner that is equally extreme, the count can never be
odd at the tail, and the smallest attainable fraction is $2/\binom{n_a+n_b}{n_a}$
rather than $1/\binom{n_a+n_b}{n_a}$. When the arms differ in size the complement
has the wrong size, is not an admissible regrouping, no mirror is forced, and
$1/\binom{n_a+n_b}{n_a}$ becomes attainable.

\section{Multiple-comparison correction}
\label{app:stats}

Correction is Holm--Bonferroni over two \emph{pre-registered} families, declared
before the final arm was trained: Family A (predictive quality, $m=4$, metric
validation task loss) and Family B (adaptive depth, $m=2$, metric depth--difficulty
correlation). The split is declared rather than chosen because the two families ask
different questions on different quantities, and pooling them would correct the
depth claim for tests that cannot bear on it. Holm controls the same familywise
error rate as plain Bonferroni and is uniformly more powerful, and here the choice
is load-bearing: plain Bonferroni over the seven comparisons would give
$0.05/7 = 0.00714$, \emph{below} the floor of 0.00794---a design in which no effect
of any size could be declared significant. Holm's easiest threshold at $m=4$ is
$0.0125$, a margin of 0.0046 above the floor.

Two consequences we hold ourselves to. A declared comparison that could not be run
still consumes its slot: dropping a failed arm must never make the survivors easier
to call significant. And a comparison whose floor exceeds its Holm threshold is
\emph{undecidable} and reported \texttt{N/A}, never ``not significant''---that
phrasing would blame the data for a limit of the design.

\section{Integrity checks and structurally-zero loss terms}
\label{app:integrity}

\subsection{Integrity checks}


Three mechanical invariants hold across the matrix, each enforced by an automated
check rather than by inspection. \textbf{(i) Inputs are bit-identical across arms
for the same record}: feature ranges are identical at $E \in \{1,5,6\}$, so a
difference between arms cannot be a difference in representation. \textbf{(ii) The
regression/classification target is absent from the input}, and no oracle routing
label is encoded into a feature slot. \textbf{(iii) The arms are parameter-matched
to 0.069\%} and the seed-blind configuration is identical within each arm, so
comparability cannot rest on a single code revision: the matrix spans four commits,
and it is the frozen specification plus the config-equality check that holds the
arms comparable, not one revision.

One disclosure belongs with the config table. Three loss terms are structurally
zero on language: the family-classification term because its head reads a detached
representation so its gradient cannot reach the trunk, and the routing-supervision
and halting-supervision terms because natural language supplies no per-token oracle
expert and no per-token oracle depth. The first is a genuine advantage over a
comparable synthetic study, where a family cross-entropy is often one of the larger
terms in the objective and every specialisation claim must be caveated for it.

\section{Loss stratified by token frequency and position}
\label{app:stratified}

\subsection{Stratified loss: the advantage concentrates on hard tokens}
\label{sec:stratified}

Fig.~\ref{fig:stratified} groups validation loss by the frequency decile of the
\emph{target} token---difficulty being a property of what must be predicted, not of
the input---and by position within the block.

Both recursive arms improve most where prediction is hardest. The MoRE$-$MoE
advantage grows from $-0.147$ nats on the most frequent decile to $-0.374$ on the
rarest; MoR's from $-0.188$ to $-0.527$. The widening is not monotone---the gap
dips at the seventh and tenth deciles---so we claim a trend, not a monotone
relationship. Neither architecture's gain is a uniform shift, which is what a
pure capacity effect would look like.

\begin{figure*}[t]
\centering
\includegraphics[width=0.98\linewidth]{fig_stratified.pdf}
\caption{(a) Validation loss by target-token frequency decile; decile 1 is the
\emph{most frequent} (easiest) and decile 10 the rarest. (b) Each recursive arm's
advantage over MoE, which widens toward the rarer deciles, though not monotonically.
Exploratory, uncorrected.}
\label{fig:stratified}
\end{figure*}

\textbf{A prediction this does not support.} If recursion were doing what the
adaptive-computation premise suggests---spending steps where context is thin---the
recursive arms should gain most at the \emph{start} of a block. They do the
opposite: MoRE's advantage over MoE is $-0.047$ nats in the first 8 positions and
$-0.359$ in the last 128; MoR's goes from $-0.250$ to $-0.465$. Both recursive arms
gain most where context is \emph{longest}. We report this because it disconfirms a
mechanism we would otherwise have asserted, and because it is consistent with the
capacity account in \S\ref{sec:mechanism}: a shared block iterated is accumulating
representation, not compensating for missing context.

\begin{figure*}[t]
\centering
\includegraphics[width=0.98\linewidth]{fig_ood.pdf}
\caption{The ordering inverts out of domain: MoRE is worst in domain (a) and best
on the Pile (b). Same three entities and colours in both panels. Exploratory,
uncorrected.}
\label{fig:ood}
\end{figure*}

\section{The subspace-trapping probes}
\label{app:trapping}

\subsection{A plausible mechanism, refuted}
\label{sec:trapping}

\textbf{[MEASURED; hypothesis falsified.]} The natural explanation for MoRE's
collapse is \emph{subspace trapping}: routing a token once confines its recursion
to a single expert's parameter subspace, and repeated application of one block
within that subspace drives the state toward a fixed point rather than refining it.
It is an appealing account and it is wrong. We tested it three ways.

\begin{center}
\begin{tabular}{lccc}
\toprule
probe & MoE & MoR & MoRE \\
\midrule
final-state participation ratio & $41.5 \pm 1.6$ & $\mathbf{70.3 \pm 5.6}$ & $57.3 \pm 2.8$ \\
between-expert mean cosine & $0.630$ & \textsc{n/a} & $0.916$ \\
\bottomrule
\end{tabular}
\end{center}

\textbf{(i) Trapping predicts lower participation ratio in the composed arm;
it is higher.} If recursion were collapsing the state onto a low-dimensional
subspace, the composed model's final-state PR should fall below MoE's, since MoE
has no recursion at all. It does not: MoRE's PR is $57.3$ against MoE's $41.5$,
and MoR's---with no routing---is higher still at $70.3$. Recursion \emph{expands}
the effective dimensionality, and routing contracts it somewhat without collapsing
it. The ordering refutes the prediction.

\textbf{(ii) Trapping predicts oscillating or fixed-point dynamics; velocity
decays smoothly.} Measuring the L2 norm of consecutive post-norm states across
depth, both recursive arms decay monotonically---MoR from 11.06 at step 2 to 5.71
at step 7, MoRE from 9.09 to 3.49---which is the signature of settling, not of
vibrating around an attractor.

\textbf{(iii) Trapping predicts that successive updates rescale the state;
they are near-orthogonal.} The mean cosine between a step's update and the state
it updates from is negative in both arms ($-0.26$ to $-0.06$), i.e. updates push
into new directions rather than amplifying existing ones. A rescaling collapse
would show positive collinearity approaching $+1$.

The recursion in this architecture is dynamically healthy: it travels, it settles,
and it adds dimensions. The failure is not a broken iterated map. That rules out
the most attractive explanation and leaves the budget competition
(\S\ref{sec:competition}) as the account the data supports---which is a stronger
position than an unexplained null, because we can say which of the obvious
mechanisms it is \emph{not}.

